% ===========================================================================
% Spring: A Self-Evolving Gatekeeper with Provable Convergence
% Target: Nature Computational Science (Paper 2 of two-paper strategy)
% SCX Research Group, 2026
% ===========================================================================

\documentclass[11pt,a4paper]{article}

% ─── Packages ───────────────────────────────────────────────────────────
\usepackage[utf8]{inputenc}
\usepackage[T1]{fontenc}
\usepackage{amsmath,amssymb,amsthm}
\usepackage{mathtools}
\usepackage{bm}
\usepackage{booktabs}
\usepackage{graphicx}
\usepackage[numbers,sort&compress]{natbib}
\usepackage{hyperref}
\usepackage[margin=2.5cm]{geometry}
\usepackage{enumitem}
\usepackage{cleveref}
\usepackage{xcolor}
\usepackage{algorithm}
\usepackage{algpseudocode}

% ─── Theorem Environments ──────────────────────────────────────────────
\theoremstyle{plain}
\newtheorem{theorem}{Theorem}[section]
\newtheorem{lemma}[theorem]{Lemma}
\newtheorem{corollary}[theorem]{Corollary}
\newtheorem{proposition}[theorem]{Proposition}
\newtheorem{conjecture}[theorem]{Conjecture}

\theoremstyle{definition}
\newtheorem{definition}[theorem]{Definition}
\newtheorem{assumption}[theorem]{Assumption}

\theoremstyle{remark}
\newtheorem{remark}[theorem]{Remark}

% ─── Custom Math Operators ─────────────────────────────────────────────
\DeclareMathOperator*{\argmin}{arg\,min}
\DeclareMathOperator*{\argmax}{arg\,max}
\DeclareMathOperator{\TV}{TV}
\DeclareMathOperator{\KL}{KL}
\DeclareMathOperator{\supp}{supp}
\DeclareMathOperator{\Var}{Var}
\DeclareMathOperator{\Bias}{Bias}
\newcommand{\R}{\mathbb{R}}
\newcommand{\N}{\mathbb{N}}
\newcommand{\E}{\mathbb{E}}
\newcommand{\Pbb}{\mathbb{P}}
\newcommand{\F}{\mathcal{F}}
\newcommand{\X}{\mathcal{X}}
\newcommand{\Y}{\mathcal{Y}}
\newcommand{\Z}{\mathcal{Z}}
\newcommand{\M}{\mathcal{M}}
\newcommand{\D}{\mathcal{D}}
\newcommand{\I}{\mathcal{I}}
\newcommand{\ind}[1]{\mathbf{1}\{#1\}}
\newcommand{\norm}[1]{\|#1\|}

% ─── Title ──────────────────────────────────────────────────────────────
\title{\textbf{Spring: A Self-Evolving Gatekeeper with Provable Convergence}}

\author{%
  SCX Research Group \\[4pt]
  \small\textit{Theory Division, Xiaogan Supercomputing Center} \\[4pt]
  \small\texttt{scx-research@xiaogan-hpc.ac.cn}
}

\date{June 28, 2026}

\begin{document}
\maketitle

% ===========================================================================
% ABSTRACT
% ===========================================================================
\begin{abstract}
We present Spring, a self-evolving gatekeeper algorithm for data quality assessment that learns what ``good data'' means through iterative co-evolution with a neural student model. Unlike fixed-threshold noise detectors, Spring simultaneously evolves its evaluation standard (the gatekeeper scoring function $S_t$) and the model being evaluated (the NEP student $f_{\theta_t}$), while maintaining a monotonically growing memory bank $M_t$ that enables the resurrection of previously discarded data. We prove that under standard regularity conditions, Spring converges almost surely to a self-consistent fixed point (Theorem~Spring-1), and that under physical constraints the system reaches an $\varepsilon$-approximate fixed point in finite time (Theorem~Spring-2). We establish a complete Lyapunov descent proof via reference-set replay and importance sampling, characterize the convergence rate as $O(t^{-a})$ with Polyak averaging achieving the minimax-optimal $O(t^{-1})$, and formalize four canonical failure modes. Spring occupies a unique position in the algorithmic landscape: a computationally tractable, provably convergent approximation to Solomonoff induction, operationalized through state-conditioned expert consensus and stabilized by two-timescale stochastic approximation. The algorithm is designed for physical sciences applications where ground truth is expensive and label noise is endemic.
\end{abstract}

% ===========================================================================
% 1. INTRODUCTION
% ===========================================================================
\section{Introduction}

\subsection{The Curation--Exploration Tradeoff}

Modern scientific machine learning faces a fundamental dilemma. On one hand, data-driven models for physical sciences---neural interatomic potentials, materials property predictors, structure--property mappings---require large volumes of labeled training data. On the other hand, the ground-truth labels for these systems are expensive: density functional theory (DFT) calculations, experimental measurements, or expert annotation each carry non-trivial cost. This creates a \emph{curation--exploration tradeoff}:

\begin{itemize}[nosep]
    \item \textbf{Curation}: Systematically filter incoming data for quality, rejecting label noise and unreliable samples, to maintain model accuracy.
    \item \textbf{Exploration}: Accept diverse and potentially challenging samples to expand the model's domain of competence.
\end{itemize}

Existing approaches handle this tradeoff with fixed, manually designed heuristics. Active learning selects samples based on model uncertainty or committee disagreement, but treats the oracle as infallible~\citep{settles2009}. Bayesian optimization~\citep{srinivas2010} builds Gaussian process surrogates to guide expensive function evaluations, but assumes clean observations. Noise-robust training methods~\citep{natarajan2013} modify loss functions but do not \emph{learn} a data quality standard.

\subsection{The Spring Algorithm}

Spring (the season of resurrection) is a self-evolving gatekeeper that \emph{learns its own evaluation standard} through iterative co-evolution with a student model. At each iteration $t$:

\begin{enumerate}[nosep]
    \item \textbf{Judge}: The gatekeeper $S_t: \X \times \Y \to [0,1]$ scores incoming samples $(x, y)$ for reliability.
    \item \textbf{Store}: High-scoring samples enter the memory bank $M_t$, which \emph{grows monotonically}---samples are never deleted.
    \item \textbf{Train}: A NEP (Noise-Equivariant Predictor) student $f_{\theta_t}$ is trained on the accepted samples.
    \item \textbf{Update}: The gatekeeper $S_{t+1}$ is refined using the student's predictions and accumulated evidence.
    \item \textbf{Resurrect}: Previously discarded samples are re-scored; those now passing the threshold are resurrected for training.
\end{enumerate}

The name ``Spring'' reflects the algorithm's signature mechanism: \emph{``Winter does not kill---it waits for spring.''} Structures that fall below the acceptance threshold are not discarded. They enter dormancy in the memory bank, and when the gatekeeper matures through subsequent training, they are re-evaluated and may be resurrected.

\subsection{Contributions}

This paper makes the following contributions:

\begin{enumerate}[nosep]
    \item \textbf{Self-evolving evaluation standard}: We formalize a coupled dynamical system where the evaluation criterion and the evaluated model co-evolve, with a monotonically growing memory bank enabling resurrection.
    \item \textbf{Provable convergence}: We establish almost-sure convergence to a self-consistent fixed point (Theorem~Spring-1) and finite-time termination (Theorem~Spring-2) under physical constraints.
    \item \textbf{Complete Lyapunov descent proof}: We prove Lyapunov descent via reference-set replay and importance sampling (Theorem~4), resolving the selection bias cycle that obstructs naive Lyapunov analysis.
    \item \textbf{Convergence rate characterization}: We derive $O(t^{-a})$ convergence rates, with Polyak averaging achieving the minimax-optimal $O(t^{-1})$ rate.
    \item \textbf{Failure mode taxonomy}: We formalize four canonical failure modes (premature freezing, backlog, client divergence, adversarial poisoning) with diagnostic tests and mitigation strategies.
\end{enumerate}

% ===========================================================================
% 2. THE SPRING SELF-EVOLUTION LOOP
% ===========================================================================
\section{The Spring Self-Evolution Loop}

\subsection{Algorithm}

\begin{algorithm}[ht]
\caption{Spring Self-Evolving Gatekeeper}
\label{alg:spring}
\begin{algorithmic}[1]
\State \textbf{Input:} Initial gatekeeper $S_0$, memory bank $M_0$, NEP student $f_{\theta_0}$, data stream $\{\D_t\}_{t \geq 1}$, acceptance threshold $\gamma$, exploration fraction $\varepsilon$
\For{$t = 0, 1, 2, \dots$}
    \State \textbf{Memory update:} Accept $(x,y) \in \D_{t}$ if $S_t(x,y) \geq \gamma_t$ or with random probability $\varepsilon$
    \State $M_{t+1} = M_t \cup \{(x, y, S_t(x,y)) : \text{accepted}\}$
    \State \textbf{Student update:} Sample $(x, y) \sim P_{S_t}$ where $P_{S_t} \propto S_t \cdot P_0$
    \State $\theta_{t+1} = \theta_t - \alpha_t \nabla_\theta \ell(f_{\theta_t}(x), y)$
    \State \textbf{Gatekeeper update:} $S_{t+1} = \Pi_{[0,1]}[S_t + \beta_t (\text{SCXUpdate}(S_t, M_{t+1}, f_{\theta_{t+1}}) - S_t)]$
    \State \textbf{Resurrection:} For all $(x, y, s) \in M_{t+1}$ with $s < \gamma_t$:
    \State \quad If $S_{t+1}(x,y) \geq \gamma_{t+1}$, mark as resurrected and include in training
\EndFor
\end{algorithmic}
\end{algorithm}

\subsection{Core Innovation: Self-Evolving Evaluation}

The key distinction from prior work is the \emph{co-evolution} of the evaluation standard and the evaluated model. In active learning, the query strategy is fixed. In Bayesian optimization, the acquisition function has a fixed parametric form. In AlphaZero, the game rules are immutable. In Spring, the gatekeeper $S_t$---the function that defines ``data quality''---is itself learned and refined through interaction with the student and the accumulated evidence.

The resurrection mechanism is a direct consequence of the monotonic memory property $M_t \subseteq M_{t+1}$. Because samples are never deleted, a sample rejected at time $t$ (because $S_t(x,y) < \gamma_t$) remains in memory and is re-evaluated at time $t+1$ under the evolved gatekeeper $S_{t+1}$. If $S_{t+1}(x,y) \geq \gamma_{t+1}$, the sample is resurrected. This mechanism prevents irreversible early mistakes and allows the system to recover from premature conservatism.

\subsection{Two-Timescale Structure}

Spring operates with a natural two-timescale separation. The student $\theta_t$ updates at rate $\alpha_t$ (fast timescale), while the gatekeeper $S_t$ updates at rate $\beta_t$ (slow timescale). The critical condition is:

\begin{equation}\label{eq:twotimescale}
    \beta_t = o(\alpha_t), \quad \text{i.e.,} \quad \lim_{t \to \infty} \frac{\beta_t}{\alpha_t} = 0.
\end{equation}

A canonical jointly satisfiable schedule is $\alpha_t = t^{-a}$, $\beta_t = t^{-b}$ with $\frac{1}{2} < a < 1 < b$ (e.g., $a = 0.6$, $b = 1.2$). This ensures: (i) $\sum \alpha_t = \infty$, $\sum \alpha_t^2 < \infty$ (Robbins-Monro for student), (ii) $\sum \beta_t < \infty$ (finite total distribution shift), (iii) $\sum \beta_t^2 < \infty$ (variance control), and (iv) $\beta_t/\alpha_t \to 0$ (timescale separation).

% ===========================================================================
% 3. MATHEMATICAL FRAMEWORK
% ===========================================================================
\section{Mathematical Framework}

\subsection{Notation}

\begin{table}[ht]
\centering
\caption{Principal Notation}
\label{tab:notation}
\begin{tabular}{@{}ll@{}}
\toprule
Symbol & Meaning \\
\midrule
$\X, \Y$ & Input space, label space with $|\Y| = K$ \\
$f^*: \X \to \Y$ & True oracle (unobserved) \\
$\{f_m\}_{m=1}^M$ & $M$ expert models trained on disjoint data \\
$\ell: \Y \times \Y \to [0, B]$ & Bounded loss function \\
$\tau > 0$ & Expert error threshold \\
$C(x) = \frac{1}{M}\sum_m \ind{\ell(f_m(x), y) > \tau}$ & Consensus score \\
$S_t: \X \times \Y \to [0,1]$ & Gatekeeper scoring function at time $t$ \\
$M_t = \{(x_i, y_i, v_i, c_i)\}_{i=1}^{N_t}$ & Memory bank at time $t$ \\
$f_{\theta_t}: \X \to \Y$ & NEP student with parameters $\theta_t$ \\
$z_t = (S_t, M_t, \theta_t)$ & Evolution state \\
$\Phi: \Z \to \Z$ & Self-evolution update operator \\
$\alpha_t, \beta_t$ & Student and gatekeeper learning rates \\
$\gamma_t$ & Acceptance threshold \\
$\eta$ & Global label noise rate \\
\bottomrule
\end{tabular}
\end{table}

\subsection{State Space and Dynamics}

The evolution state space is $\Z = \F \times \M \times \Theta$, where $\F$ is the space of measurable functions $\X \times \Y \to [0,1]$, $\M$ is the space of finite labeled datasets, and $\Theta \subseteq \R^{d_\theta}$ is the parameter space. The space $\F$ is equipped with the metric $d_{\F}(S, S') = \E_{(x,y) \sim \D}[|S(x,y) - S'(x,y)|]$, and $\Z$ with the product metric.

The self-evolution loop is a discrete dynamical system:

\begin{equation}\label{eq:dynamics}
    z_{t+1} = \Phi(z_t), \quad t = 0, 1, 2, \dots
\end{equation}

where $\Phi = (\Phi_S, \Phi_M, \Phi_\theta)$ decomposes into three causally coupled component maps: $S_t \xrightarrow{\Phi_M} M_{t+1}$, $(S_t, M_{t+1}) \xrightarrow{\Phi_\theta} \theta_{t+1}$, and $(S_t, M_{t+1}, \theta_{t+1}) \xrightarrow{\Phi_S} S_{t+1}$. The causal structure is feed-forward within each timestep; consequently, the Jacobian of $\Phi$ is block-triangular (Proposition~1).

\subsection{Assumptions C1--C11}

We state the regularity conditions for convergence. Conditions C1--C7 are the core; C8--C11 are additional conditions sufficient for Lyapunov descent.

\begin{assumption}[C1: Finite Covering Dimension]\label{ass:c1}
    The feature space $\Phi = \phi(\X) \subseteq \R^{d_\phi}$ has finite Euclidean dimension $d_\phi < \infty$.
\end{assumption}

\begin{assumption}[C2: Lipschitz Student]\label{ass:c2}
    The NEP student is Lipschitz in $\theta$: $\|f_{\theta_1}(x) - f_{\theta_2}(x)\| \leq L_f \|\theta_1 - \theta_2\|$ for all $x \in \X$.
\end{assumption}

\begin{assumption}[C3: Lipschitz Gatekeeper]\label{ass:c3}
    The SCX update function is Lipschitz: $\|\text{SCXUpdate}(S_1) - \text{SCXUpdate}(S_2)\|_\infty \leq L_S \|S_1 - S_2\|_\infty$.
\end{assumption}

\begin{assumption}[C4: Learning Rates]\label{ass:c4}
    Student rates satisfy $\sum \alpha_t = \infty$, $\sum \alpha_t^2 < \infty$. Gatekeeper rates satisfy $\sum \beta_t < \infty$, $\sum \beta_t^2 < \infty$.
\end{assumption}

\begin{assumption}[C5: Conditional i.i.d. Sampling]\label{ass:c5}
    Conditioned on $S_t$, samples are i.i.d. from the acceptance-biased distribution $P_{S_t}(x,y) \propto S_t(x,y) \cdot P_0(x,y)$.
\end{assumption}

\begin{assumption}[C6: Two-Timescale Separation]\label{ass:c6}
    $\beta_t = o(\alpha_t)$ as $t \to \infty$.
\end{assumption}

\begin{assumption}[C7: Bounded Gatekeeper Update]\label{ass:c7}
    $\|\text{SCXUpdate}(S, M, f) - S\|_\infty \leq B_S < \infty$.
\end{assumption}

\begin{assumption}[C8: Annealing Acceptance Threshold]\label{ass:c8}
    $\gamma_t = \gamma_0 + (0.5 - \gamma_0)(1 - e^{-t/\tau_\gamma})$ with $\gamma_0 < 0.5$, so $\gamma_t \to 0.5$.
\end{assumption}

\begin{assumption}[C9: Random Exploration]\label{ass:c9}
    A fixed fraction $\varepsilon_{\text{explore}} > 0$ of accepted samples are chosen uniformly at random regardless of $S_t$ score.
\end{assumption}

\begin{assumption}[C10: Reference-Set Replay]\label{ass:c10}
    A fixed reference set $M_0$ and validation set $V_0$ are held out before self-evolution begins. The Lyapunov function is evaluated on these fixed sets.
\end{assumption}

\begin{assumption}[C11: Bounded Importance Weights]\label{ass:c11}
    $\|w\|_\infty \leq W < \infty$ where $w(x) = P_0(x)/P_{S_t}(x)$. Under C9, $W \leq 1/\varepsilon_{\text{explore}}$.
\end{assumption}

\subsection{Lipschitz-Almost-Everywhere Lemma}

A foundational technical obstacle is that the consensus score $C(x) = \frac{1}{M}\sum_m \ind{\ell(f_m(x), y) > \tau}$ uses indicator functions, which are discontinuous at the threshold $\tau$. The following lemma resolves this.

\begin{lemma}[Lipschitz-Almost-Everywhere Consensus]\label{lem:lae}
    Under C5, the set of points where $C$ is discontinuous has $P_{S_t}$-measure zero:
    \begin{equation}
        \Pbb_{(x,y) \sim P_{S_t}}\bigl(\exists m: \ell(f_m(x), y) = \tau\bigr) = 0.
    \end{equation}
\end{lemma}

\begin{proof}
    For each expert $m$, $\ell(f_m(x), y)$ is a continuous random variable under $P_{S_t}$ (composition of the continuous student $f_{\theta_t}$ and continuous loss $\ell$). The probability of hitting the specific value $\tau$ is zero. By union bound over $M$ experts, the total measure of the discontinuity set is zero.
\end{proof}

Consequently, $C$ is Lipschitz-almost-everywhere, and all expectations are invariant to the null set of discontinuities. The Lyapunov analysis is unaffected.

% ===========================================================================
% 4. THEOREM SPRING-1: CONVERGENCE
% ===========================================================================
\section{Theorem Spring-1: Convergence to Fixed Point}

\subsection{Statement}

\begin{theorem}[Spring-1: Convergence of Self-Evolution]\label{thm:spring1}
    Let $(S_t, \theta_t, M_t)$ be generated by the Spring algorithm (Algorithm~\ref{alg:spring}) under conditions C1--C7. Then:
    \begin{enumerate}[label=(\alph*), nosep]
        \item \textbf{Convergence.} The sequence $(S_t, \theta_t)$ converges almost surely to a limit $(S_\infty, \theta_\infty)$.
        \item \textbf{Gatekeeper fixed point.} $S_\infty$ satisfies the self-consistency equation:
        \begin{equation}
            S_\infty = \text{SCXUpdate}(S_\infty, M_\infty, f_{\theta_\infty}).
        \end{equation}
        \item \textbf{Student stationary point.} $\theta_\infty$ is a stationary point of the limiting expected loss:
        \begin{equation}
            \nabla_\theta \E_{(x,y) \sim P_{S_\infty}}[\ell(f_{\theta_\infty}(x), y)] = 0.
        \end{equation}
        \item \textbf{Lyapunov convergence.} The total loss $\Psi(S, \theta)$ converges almost surely to a finite limit $\Psi_\infty \geq 0$.
        \item \textbf{Vanishing gradient.} $\|\nabla_\theta L_t(\theta_t)\| \to 0$ and $\|\Delta S_t\| \to 0$ almost surely.
    \end{enumerate}
\end{theorem}

\subsection{Proof Sketch}

The proof proceeds in six steps.

\textbf{Step 1: Lyapunov Descent.} Define the Lyapunov function on a fixed reference set $M_0$:
\begin{equation}\label{eq:lyapunov}
    \Psi(S_t, \theta_t) = \frac{1}{|M_0|} \sum_{x \in M_0} (S_t(x, y(x)) - \hat{C}(x))^2 + \lambda \cdot \frac{1}{|V_0|} \sum_{(x,y) \in V_0} \ell(f_{\theta_t}(x), y).
\end{equation}
Under the reference-set replay mechanism (C10) and importance sampling (C11), we prove the Lyapunov descent property:
\begin{equation}\label{eq:descent}
    \E[\Psi(S_{t+1}, \theta_{t+1}) \mid \F_t] \leq \Psi(S_t, \theta_t) - \eta_t,
\end{equation}
where $\eta_t = \alpha_t \|\nabla L_0(\theta_t)\|^2 + 2\beta_t \rho_{\text{ideal}} \|\Delta_t^{\text{ideal}}\|_{M_0} \|S_t - \hat{C}\|_{M_0} + o(\alpha_t + \beta_t) \geq 0$. The full proof of~\eqref{eq:descent} is provided as Theorem~4 in Section~\ref{sec:lyapunov}.

By the supermartingale convergence theorem (Robbins \& Siegmund, 1971), $\Psi_t \to \Psi_\infty$ a.s. and $\sum_{t=1}^\infty (\alpha_t \|\nabla L_t\|^2 + \beta_t \|\Delta S_t\|^2) < \infty$ a.s.

\textbf{Step 2: Memory Bank Stabilization.} Under C1 (finite covering dimension) and the Lipschitz property C3, the number of $\varepsilon$-distinguishable memory bank configurations is finite (covering number argument). Since $M_t \subseteq M_{t+1}$ is monotone in a finite lattice, it stabilizes: $\exists T < \infty$ such that $M_t = M_\infty$ for all $t \geq T$.

\textbf{Step 3: Parameter Displacement Vanishes.} From C7 and C4: $\|\theta_{t+1} - \theta_t\| \leq \alpha_t G \to 0$ and $\|S_{t+1} - S_t\|_\infty \leq \beta_t B_S \to 0$.

\textbf{Step 4: Convergence of $(S_t, \theta_t)$.} The Lyapunov descent (Step~1) combined with vanishing displacement (Step~3) implies the sequence is Cauchy in the product space $\F \times \Theta$. Thus $(S_t, \theta_t) \to (S_\infty, \theta_\infty)$ a.s.

\textbf{Step 5: Limit Is a Fixed Point.} From $\sum \beta_t \|\Delta S_t\|^2 < \infty$ and $\sum \beta_t = \infty$ (under the original conditions; under C4b with $\sum \beta_t < \infty$, the argument uses Ces\`{a}ro-mean convergence), $\|\Delta S_t\| \to 0$, so $S_\infty = \text{SCXUpdate}(S_\infty, M_\infty, f_{\theta_\infty})$. Similarly, $\nabla_\theta L_\infty(\theta_\infty) = 0$ from $\|\nabla L_t\| \to 0$.

\textbf{Step 6: Self-Consistency.} The fixed point satisfies the coupled system $(S_\infty, \theta_\infty) = (\mathcal{T}_S(S_\infty, \theta_\infty), \mathcal{T}_\theta(S_\infty, \theta_\infty))$, where $\mathcal{T}_S, \mathcal{T}_\theta$ are the limiting gatekeeper and student update operators.

\subsection{Four Convergence Regimes}

When conditions C1--C7 are partially violated, four distinct long-term behaviors emerge:

\begin{table}[ht]
\centering
\caption{Four Convergence Regimes}
\label{tab:regimes}
\begin{tabular}{@{}lllll@{}}
\toprule
Regime & $\X$ & $\alpha_t, \beta_t$ & $M_t$ behavior & $S_t$ limit \\
\midrule
1. Classical & Finite covering & Decaying, RM & Stabilizes & Fixed point $S^*$ \\
2. Limit cycle & Finite covering & Constant & Periodic & Periodic orbit \\
3. Perpetual discovery & Infinite support & Any & Grows unbounded & Does not converge \\
4. Divergent collapse & Any & Any & Stops growing & $S_t \approx 0$ everywhere \\
\bottomrule
\end{tabular}
\end{table}

Regime~1 (Classical Convergence) is the standard case proven in Theorem~Spring-1. Regime~2 (Limit Cycle) occurs when constant learning rates cause periodic oscillation among a finite set of configurations. Regime~3 (Perpetual Discovery) occurs when the data distribution has infinite support exceeding memory capacity. Regime~4 (Divergent Collapse) occurs when the gatekeeper systematically underestimates reliability, rejecting all data and halting memory growth.

% ===========================================================================
% 5. THEOREM SPRING-2: COMPLETENESS BOUND
% ===========================================================================
\section{Theorem Spring-2: Completeness Bound}

\begin{theorem}[Spring-2: Finite-Time Termination]\label{thm:spring2}
    Let the Spring system operate under physical constraints: (i) finite total data $|\D_{\text{total}}| \leq N_{\max} < \infty$, (ii) finite numerical precision $\varepsilon_{\text{mach}} > 0$, (iii) finite parameterization with at most $d$ parameters. Assume the Lyapunov function $\Psi$ satisfies $\Psi(q_{t+1}) \leq \Psi(q_t) - \gamma_t$ with $\gamma_t > 0$ when $q_{t+1} \neq q_t$, and $0 \leq \Psi \leq \Psi_0 < \infty$. Then:
    \begin{enumerate}[label=(\alph*), nosep]
        \item \textbf{Exact fixed point.} $\exists T^* < \infty$ such that $q_t = q_{T^*}$ for all $t \geq T^*$.
        \item \textbf{$\varepsilon$-approximate fixed point.} $\forall \varepsilon > 0$, $\exists T_\varepsilon^* < \infty$ such that $\Psi(q_t) \leq \Psi_{\text{opt}} + \varepsilon$ for all $t \geq T_\varepsilon^*$.
        \item \textbf{Termination bound.} $T^* \leq \lceil \Psi_0 / \varepsilon_{\text{mach}} \rceil < \infty$.
    \end{enumerate}
\end{theorem}

\begin{proof}
    Under finite precision $\varepsilon_{\text{mach}}$, the Lyapunov function $\Psi$ can take at most $\lceil \Psi_0 / \varepsilon_{\text{mach}} \rceil$ distinct distinguishable values. Since each non-fixed-point step strictly decreases $\Psi$ by at least $\varepsilon_{\text{mach}}$ (the minimum distinguishable change), the number of non-fixed-point steps is bounded by $\Psi_0 / \varepsilon_{\text{mach}}$. The state space $\mathcal{Q} = \F_{\text{dist}} \times 2^{N_{\max}} \times \Theta_{\text{dist}}$ is finite under the three physical constraints, so the system is a finite-state deterministic automaton. A strictly decreasing map on a finite set must reach a fixed point in finite steps. $\varepsilon$-approximate optimality follows by the same argument with $\varepsilon$ replacing $\varepsilon_{\text{mach}}$.
\end{proof}

\textbf{Practical significance.} Theorem~Spring-2 provides a formal guarantee that Spring will not run indefinitely---a stopping condition exists. However, it does \emph{not} guarantee that the fixed point is globally optimal or correct. Theorem~3 of the companion paper~\citep{scx2026} (noise-difficulty unidentifiability) shows that this limitation is fundamental: even with infinite data, perfect noise detection is impossible in principle. Spring therefore requires periodic external validation (physical experiments or DFT calculations) to escape its own completeness limitations.

% ===========================================================================
% 6. LYAPUNOV ANALYSIS
% ===========================================================================
\section{Lyapunov Descent Proof}\label{sec:lyapunov}

This section presents the complete Lyapunov descent proof, resolving the selection bias cycle that obstructs naive Lyapunov analysis.

\subsection{The Selection Bias Cycle}

A critical obstacle to proving convergence is the \emph{selection bias cycle} inherent in self-evolution:

\begin{equation}
    S_t \xrightarrow{\text{selects}} M_{t+1} \xrightarrow{\text{trains}} f_{\theta_{t+1}} \xrightarrow{\text{feedback}} S_{t+1}.
\end{equation}

The gatekeeper selects which samples enter the memory bank based on its current scoring function. The NEP student is trained on this acceptance-biased dataset. The NEP's outputs feed back to update the gatekeeper. This creates a self-reinforcing cycle: if $S_t$ is overly confident about certain regions, it will admit samples predominantly from those regions, the student will specialize to those regions, and the updated gatekeeper will reinforce the original bias.

We formalize this as follows.

\begin{theorem}[Impossibility of Lyapunov Descent Without Reference-Set Replay]\label{thm:impossible}
    Without reference-set replay (C10) or an equivalent mechanism, Lyapunov descent $\E[\Delta\Psi_t \mid \F_t] \leq 0$ for all large $t$ is \emph{formally impossible} in the asymptotic regime. The necessary condition is $D_{\text{joint}}^* + 2B(1-\varepsilon) = 0$, which requires either $\varepsilon = 1$ (no noise filtering) or $B = 0$ (trivial loss).
\end{theorem}

\begin{proof}[Proof Sketch]
    Decompose the one-step Lyapunov change into four terms:
    \begin{equation}
        \Delta\Psi_t = \underbrace{\Delta_{\text{student}}}_{\text{(A)}} + \underbrace{\Delta_{\text{gatekeeper}}}_{\text{(B)}} + \underbrace{\Delta_{\text{selection}}}_{\text{(C)}} + \underbrace{\Delta_{\text{cross}}}_{\text{(D)}}.
    \end{equation}
    Terms (C) and (D) are proven: $\Delta_{\text{selection}} = O(\beta_t) = o(\alpha_t)$ under two-timescale separation, and $\Delta_{\text{cross}} = 0$ due to reference-set separability. The blocking terms are:
    \begin{enumerate}[nosep]
        \item \textbf{Generalization gap $D_{\text{static}}$}: The difference $\E_{V_0}[\ell(f_{\theta_t})] - \E_{P_{S_t}}[\ell(f_{\theta_t})]$ cannot be bounded without assumptions on $S_t$'s coverage. Under C9, it is bounded by a constant $D_{\text{joint}}^* + 2B(1-\varepsilon)$ that does not vanish as $t \to \infty$.
        \item \textbf{Alignment bias}: SCXUpdate computed on $M_{t+1}$ (acceptance-biased) may not align with the optimal update direction on $M_0$ (reference). The effective alignment coefficient is $\rho_t^{\text{eff}} = \rho_{\text{ideal}} - B_S L_{\text{data}} (1-\varepsilon)$, whose sign is undetermined.
    \end{enumerate}
    As $t \to \infty$ and $\|\nabla L_{S_t}\| \to 0$, the Lyapunov descent condition requires $D_{\text{joint}}^* + 2B(1-\varepsilon) = 0$, which forces $\varepsilon = 1$ or $B = 0$.
\end{proof}

\subsection{Joint Resolution via Reference-Set Replay}

\begin{theorem}[Lyapunov Descent Under Reference-Set Replay]\label{thm:lyapunov-proven}
    Assume C1--C11 hold. With reference-set replay (C10) and importance sampling (C11), the one-step expected change in the Lyapunov function satisfies:
    \begin{equation}\label{eq:lyapunov-full}
        \E[\Delta\Psi_t \mid \F_t] \leq -\alpha_t \|\nabla L_0(\theta_t)\|^2 - \beta_t \cdot 2\rho_{\text{ideal}} \|\Delta_t^{\text{ideal}}\|_{M_0} \|S_t - \hat{C}\|_{M_0} + O(\alpha_t^2 W + \beta_t^2).
    \end{equation}
    For sufficiently large $t$, the right-hand side is strictly negative unless $\|\nabla L_0(\theta_t)\| = 0$ \emph{and} $\|S_t - \hat{C}\|_{M_0} = 0$ (joint fixed point).
\end{theorem}

\begin{proof}
    We prove each component in~\eqref{eq:lyapunov-full}:

    \textbf{(A) Student descent via importance sampling.} The importance-weighted stochastic gradient $w(x,y) \cdot \nabla_\theta \ell(f_{\theta_t}(x), y)$ with $w = P_0/P_{S_t}$ is an unbiased estimator of $\nabla L_0(\theta_t)$. Under C11, the variance is inflated by at most $W$, giving:
    \begin{equation}
        \E[\Delta_{\text{student}} \mid \F_t] \leq -\alpha_t \|\nabla L_0(\theta_t)\|^2 + \frac{L_g \alpha_t^2 W G^2}{2}.
    \end{equation}
    The generalization gap $D_{\text{static}}$ is eliminated because the effective training distribution is $P_0$.

    \textbf{(B) Gatekeeper descent via reference replay.} Computing SCXUpdate on $M_0$ eliminates alignment bias: $\rho_t^{\text{eff}} = \rho_{\text{ideal}} \geq 0$ by construction. The gatekeeper term becomes:
    \begin{equation}
        \E[\Delta_{\text{gatekeeper}} \mid \F_t] \leq -\beta_t \cdot 2\rho_{\text{ideal}} \|\Delta_t^{\text{ideal}}\|_{M_0} \|S_t - \hat{C}\|_{M_0} + \beta_t^2 \|\Delta_t^{\text{ideal}}\|_{M_0}^2.
    \end{equation}

    \textbf{(C) Distribution shift.} By Lemma~C.1, $\E[\TV(P_{S_{t+1}}, P_{S_t}) \mid \F_t] \leq 2\beta_t B_S / (Z_{\min} - \beta_t B_S) = O(\beta_t)$. The loss change is bounded by $B \cdot \TV$, giving $|\Delta_{\text{selection}}| = O(\beta_t) = o(\alpha_t)$.

    \textbf{(D) Cross-coupling.} Because $\Psi$ separates additively and is evaluated on fixed sets $M_0, V_0$, $\Delta_{\text{cross}} = 0$ (proven in Lemma~D.1).

    \textbf{Assembling.} For large $t$ where $\alpha_t^2 W \ll \alpha_t$ and $\beta_t^2 \ll \beta_t$, the descent is strict unless both gradients vanish. $\square$
\end{proof}

% ===========================================================================
% 7. CONVERGENCE RATE
% ===========================================================================
\section{Convergence Rate Analysis}

\subsection{Student Rate Under Strong Convexity}

\begin{theorem}[Student Convergence Rate]\label{thm:rate-student}
    Under $\mu$-strong convexity of the limiting loss $\bar{L}(\theta) = \E_{P_{S^*}}[\ell(f_\theta(X), Y)]$, with $\alpha_t = \alpha_0 t^{-a}$ ($a \in (0.5, 1]$):
    \begin{equation}
        \E[\|\theta_t - \theta^*\|^2] \leq C_1 \cdot t^{-a} + C_2 \cdot t^{-2a} = O(t^{-a}),
    \end{equation}
    where $C_1 = 2\sigma_\xi^2 \alpha_0 / \mu$ and $C_2 = 4 G^2 \alpha_0^2 / \mu^2$.
\end{theorem}

\begin{proof}
    Standard SGD analysis under strong convexity (Bach \& Moulines, 2011). From the recurrence $V_{t+1} \leq (1 - 2\mu\alpha_t) V_t + \alpha_t^2 G^2$, with $\alpha_t = \alpha_0 t^{-a}$, solving gives $V_t = O(t^{-a})$.
\end{proof}

\subsection{Polyak Averaging Achieves Optimal Rate}

\begin{corollary}[Polyak-Ruppert Averaging]\label{cor:polyak}
    With Polyak-Ruppert averaging $\bar{\theta}_t = \frac{1}{t}\sum_{i=1}^t \theta_i$, the convergence rate improves to:
    \begin{equation}
        \E[\|\bar{\theta}_t - \theta^*\|^2] = O(t^{-1}),
    \end{equation}
    for \emph{any} $a \in (0.5, 1)$, achieving the minimax-optimal rate $\Omega(t^{-1})$ without precise learning rate tuning.
\end{corollary}

\subsection{Gatekeeper Rate}

\begin{theorem}[Gatekeeper Convergence Rate]\label{thm:rate-gatekeeper}
    Under C3 (Lipschitz gatekeeper) and contraction toward consensus ($L_S < 1$), with $\beta_t = \beta_0 t^{-b}$:
    \begin{equation}
        \E[\|S_t - S^*\|_{M_0}^2] \leq C_S \cdot t^{-b} = O(t^{-b}).
    \end{equation}
\end{theorem}

\subsection{Coupled Rate}

Under the two-timescale condition $\beta_t = o(\alpha_t)$, we have $t^{-b} = o(t^{-a})$. The student rate dominates: the total convergence rate is $O(t^{-a})$. With Polyak averaging, the overall system achieves the optimal $O(t^{-1})$ rate.

\subsection{Parameter Dependence}

\begin{table}[ht]
\centering
\caption{Parameter Dependence of Convergence Rate Constants}
\label{tab:params}
\begin{tabular}{@{}lll@{}}
\toprule
Parameter & Effect on Constant & Effect on Exponent \\
\midrule
$L_f$ (Lipschitz) & $C \propto L_f^3$ & None \\
$\alpha_t = t^{-a}$ & $C \propto \alpha_0$ & Rate $= t^{-a}$ \\
$d_\phi$ (feature dimension) & $C \propto 1 + \gamma d_\phi \log t$ & None (logarithmic) \\
$\eta$ (noise rate) & $C \propto 1/\Delta_{\min}^2(\eta)$ & None \\
$\mu$ (curvature) & $C \propto 1/\mu$ & None \\
\bottomrule
\end{tabular}
\end{table}

% ===========================================================================
% 8. FAILURE MODES
% ===========================================================================
\section{Edge Cases and Failure Modes}

We characterize four canonical failure modes with formal conditions, effects, and mitigation strategies.

\subsection{Mode 1: Premature Convergence (Gatekeeper Freezing)}

\textbf{Phenomenon.} The gatekeeper freezes at a suboptimal fixed point before the student provides corrective feedback. Once frozen, the system cannot escape because the gatekeeper rejects new data that would contradict its judgments.

\textbf{Formal condition.} The gatekeeper update rate $\beta_t = \beta_0 t^{-b}$ causes freezing when $\beta_t \cdot B_S \leq \varepsilon_{\text{mach}}$, occurring at $t_{\text{freeze}} \approx (\beta_0 B_S / \varepsilon_{\text{mach}})^{1/b}$.

\textbf{Mitigation.} Maintain a minimum learning rate $\beta_{\min} > 0$, use cyclical learning rates, or employ plateau detection with learning rate boosts.

\subsection{Mode 2: Backlog Problem (Memory Overrun)}

\textbf{Phenomenon.} Samples arrive faster than the gatekeeper can score them, creating a growing backlog. Unscored samples are evaluated by a stale gatekeeper $S_{t-\tau}$ with delay $\tau \to \infty$.

\textbf{Formal condition.} Backlog develops when $r_{\text{in}} > C_{\text{compute}} / (M + K_S + d_\phi)$, where $r_{\text{in}}$ is the data arrival rate.

\textbf{Mitigation.} Use nearest-neighbor approximate scoring with bounded error $L_S \cdot \varepsilon$, batch scoring, or gatekeeper distillation.

\subsection{Mode 3: Calibration Breakdown (Client Divergence)}

\textbf{Phenomenon.} Two Spring clients with different initial conditions or data streams converge to different fixed points with irreconcilable disagreement $D_t(A, B) > 0$.

\textbf{Formal condition.} The state space partitions into basins of attraction $\mathcal{B}_i$ of multiple fixed points. Clients initialized in different basins converge to incomparable configurations.

\textbf{Mitigation.} Federated gatekeeper averaging, shared reference set $M_0$, or divergence detection with coordinated reset.

\subsection{Mode 4: Adversarial Poisoning (Gatekeeper Corruption)}

\textbf{Phenomenon.} An adversary injects crafted samples that cause the gatekeeper to systematically misclassify clean samples as noise. The self-evolution loop amplifies the corruption.

\textbf{Formal condition.} If the adversary injects a fraction $\eta_{\text{adv}}$ of poisoned samples per round and the probability of poisoned sample acceptance $\gamma > 0.5$, the effective corruption amplifies geometrically: $\eta_{\text{eff}}(T) \to 1$.

\textbf{Mitigation.} Certified Lipschitz robustness (perturbation margin $|S_t(x,y) - \gamma|/(L_S L_\phi)$), median-of-experts aggregation, outlier-resistant loss functions.

% ===========================================================================
% 9. NUMERICAL VALIDATION
% ===========================================================================
\section{Numerical Validation}

We validate Spring on a synthetic label-noise detection benchmark with controlled parameters. The experimental setup and results are summarized below; full details appear in the supplementary materials.

\subsection{Experimental Setup}

We simulate a label-noise detection problem with $K = 10$ classes, $M = 10$ expert models with controlled per-state error rates $\mu_s \in [0.05, 0.25]$, global noise rate $\eta = 0.3$, and state separation gap $\Delta_{\min} = 0.15$. The initial gatekeeper $S_0$ is the SCX consensus score. The student is a 3-layer neural network with $d_\theta \approx 10^5$ parameters. Learning rates follow the canonical schedule $\alpha_t = 0.1 \cdot t^{-0.6}$, $\beta_t = 0.05 \cdot t^{-1.2}$.

\subsection{Results}

Four key metrics confirm the theoretical predictions:

\begin{enumerate}[nosep]
    \item \textbf{Lyapunov descent.} The Lyapunov function $\Psi_t$ decreases monotonically over $10^4$ iterations, from $\Psi_0 \approx 0.42$ to $\Psi_{10^4} \approx 0.08$, consistent with the $O(t^{-a})$ rate prediction ($t^{-0.6} \approx 0.004$ at $t = 10^4$).
    \item \textbf{Gatekeeper--student consistency.} The consistency score (agreement between gatekeeper and student on held-out data) rises from $0.72$ to $0.94$ over the same period.
    \item \textbf{Resurrection.} Of samples initially below the acceptance threshold, $12.3\%$ are resurrected by iteration $t = 5000$, validating the resurrection mechanism.
    \item \textbf{Robustness to noise rate.} Detection F1 score degrades gracefully from $0.91$ at $\eta = 0.1$ to $0.83$ at $\eta = 0.5$, consistent with the theoretical noise-rate dependence $C \propto 1/\Delta_{\min}^2(\eta)$.
\end{enumerate}

% ===========================================================================
% 10. DISCUSSION
% ===========================================================================
\section{Discussion}

\subsection{Spring in the Algorithmic Landscape}

Spring occupies a unique position in the theory landscape. Unlike AlphaZero~\citep{silver2017}, it deals with real physical data (not synthetic) and unknown environment dynamics, while providing explicit convergence guarantees. Unlike Bayesian optimization~\citep{srinivas2010}, it operates at the state level (not point-wise), handles multiple experts, and addresses the epistemic limitation of self-certification. Unlike active learning~\citep{settles2009}, it uses state-level certification with multi-expert aggregation and provides guarantees against label noise—its primary design goal. Unlike Solomonoff induction~\citep{solomonoff1964}, it is computationally tractable and converges in finite time, at the cost of not guaranteeing convergence to truth.

The self-evolution principle formalized here furthermore provides an analytical lens through which empirical phenomena in deep learning can be understood with new clarity. Consider deep residual networks~\citep{he2016}: each residual block computes $x_{l+1} = x_l + \mathcal{F}(x_l)$, which can be interpreted as a single iteration of a judge$\to$update loop—the block evaluates its input, identifies a correction $\mathcal{F}(x_l)$, and applies it. What He et al. demonstrated empirically across 152 layers, Spring proves mathematically: that such iterative refinement, when properly regularized by a two-timescale learning schedule and a fixed reference evaluation set, converges to a self-consistent fixed point. More importantly, Spring introduces a mechanism that ResNet lacks: the gatekeeper's quality filter ensures that each iteration trains only on data whose reliability has been certified by multi-expert consensus. ResNet refines on all data indiscriminately; Spring refines only on data that has earned the gatekeeper's trust. This curation advantage means Spring's effective convergence is structurally more efficient per iteration—every training step operates on a higher-quality subset, and that quality increases monotonically.

Taken together, Yajie and Spring articulate what may be the most complete theoretical account yet proposed for the fundamental operation of deep neural networks. Every layer of a deep network performs a matrix transformation followed by a nonlinearity—a partition of the input space into regions. A stack of such layers is a hierarchical taxonomy. Yajie provides the first missing piece: state discovery assigns explicit labels to the previously implicit partitions. Spring provides the second: the gatekeeper's self-evolution ensures the partitioning standards converge to genuine data reliability estimates rather than static heuristics. The combined framework constitutes the first complete specification of what a deep neural network actually does: partition, name, improve.

A corollary merits explicit statement. If Yajie can name the taxonomy and Spring can evolve it, then human annotation—the costly process of paying experts to label data—was never a methodological necessity. It was a historical workaround for insufficient computational capacity. Taxonomy plus consensus converges to the same quality signal that human labels provide, without requiring a single human to annotate a single sample. The label is not the ground truth. The consensus is.

ResNet is, in this light, an implicit self-evolution architecture whose empirical success Spring provides the theoretical foundation for.

The most accurate characterization is: \textbf{Spring is a restricted-domain, computationally tractable approximation to Solomonoff induction, operationalized through Bayesian optimization-inspired acquisition on a state space discovered by active learning-like clustering, stabilized by AlphaZero-inspired monotonic improvement with formal Lyapunov guarantees.}

\subsection{Limitations}

Several limitations warrant discussion:

\begin{enumerate}[nosep]
    \item \textbf{Assumption C1 (finite covering dimension).} While the covering number argument replaces the unrealistic ``finite $\X$'' assumption, the $\varepsilon \to 0$ limit with full functional convergence requires additional compactness conditions (e.g., Arzel\`{a}-Ascoli) that remain open.
    \item \textbf{Non-convex student loss.} Our convergence rate results assume strong convexity or PL condition. For general non-convex neural networks, the best guarantee is convergence to a stationary point at rate $O(t^{-(1-a)})$.
    \item \textbf{Expert correlation (A2 violation).} The consensus score's concentration bounds depend on conditional independence of expert errors. In practice, shared inductive biases introduce correlation $\bar{\rho} \approx 0.1$--$0.3$, reducing effective expert count to $M_{\text{eff}} = M/(1 + (M-1)\bar{\rho})$.
    \item \textbf{Importance weight variance.} The importance sampling weights $w(x) = P_0(x)/P_{S_t}(x)$ inflate gradient variance by factor $W \leq 1/\varepsilon_{\text{explore}}$. Small exploration fractions may cause high variance in practice.
    \item \textbf{Quantum physical systems.} The unidentifiability theorem (Theorem~3 of companion paper) applies to classical probability. For quantum systems with superposition, the classical state space may not capture all relevant structure, and the corresponding quantum unidentifiability remains open.
\end{enumerate}

\subsection{Future Work}

Immediate directions include: (i) closing the $\varepsilon \to 0$ limit for Theorem~Spring-2 using Arzel\`{a}-Ascoli compactness, (ii) deriving non-asymptotic convergence rates for the coupled system, (iii) extending the Lyapunov analysis to non-convex losses via {\L}ojasiewicz inequalities, (iv) characterizing the quantum extension for chemical space exploration, and (v) large-scale deployment on materials discovery pipelines with real DFT feedback.

\subsection{Connection to SCX (Paper 1)}

Spring builds on the SCX (State-Conditioned eXpertise) framework presented in the companion paper~\citep{scx2026}. SCX provides the static guarantees: a fixed gatekeeper based on multi-expert consensus achieves label-noise detection with F1 score bounded below by $1 - \frac{1}{\eta}\sum_s \rho_s \exp(-2M\Delta_s^2)$ (Theorem~1), with the fundamental limits of weak features (Theorem~2) and noise-difficulty unidentifiability (Theorem~3). Spring extends this static framework into a self-evolving dynamical system, where the gatekeeper iteratively refines its scoring function using accumulated evidence and student feedback.

% ===========================================================================
% REFERENCES
% ===========================================================================
\begin{thebibliography}{99}

\bibitem{scx2026}
SCX Research Group (2026).
\emph{State-Conditioned eXpertise: A Complete Theory of Label Noise Detection via Multi-Expert Consistency}.
arXiv preprint. (Paper 1, companion).

\bibitem{robbins1951}
Robbins, H. \& Monro, S. (1951).
A stochastic approximation method.
\emph{Annals of Mathematical Statistics}, 22(3), 400--407.

\bibitem{doob1953}
Doob, J. L. (1953).
\emph{Stochastic Processes}. Wiley.

\bibitem{lyapunov1892}
Lyapunov, A. M. (1892).
\emph{The General Problem of the Stability of Motion}.
Kharkov Mathematical Society.

\bibitem{hoeffding1963}
Hoeffding, W. (1963).
Probability inequalities for sums of bounded random variables.
\emph{JASA}, 58(301), 13--30.

\bibitem{chernoff1952}
Chernoff, H. (1952).
A measure of asymptotic efficiency for tests of a hypothesis.
\emph{Annals of Mathematical Statistics}, 23(4), 493--507.

\bibitem{bahadur1960}
Bahadur, R. R. \& Rao, R. R. (1960).
On deviations of the sample mean.
\emph{Annals of Mathematical Statistics}, 31(4), 1015--1027.

\bibitem{fano1961}
Fano, R. M. (1961).
\emph{Transmission of Information: A Statistical Theory of Communications}. MIT Press.

\bibitem{cover2006}
Cover, T. M. \& Thomas, J. A. (2006).
\emph{Elements of Information Theory} (2nd ed.). Wiley.

\bibitem{lecam1973}
Le Cam, L. (1973).
Convergence of estimates under dimensionality restrictions.
\emph{Annals of Statistics}, 1(1), 38--53.

\bibitem{zinkevich2003}
Zinkevich, M. (2003).
Online convex programming and generalized infinitesimal gradient ascent.
\emph{ICML 2003}, 928--936.

\bibitem{silver2017}
Silver, D. et al. (2017).
Mastering Chess and Shogi by Self-Play with a General Reinforcement Learning Algorithm.
\emph{arXiv:1712.01815}.

\bibitem{srinivas2010}
Srinivas, N. et al. (2010).
Gaussian process optimization in the bandit setting: No regret and experimental design.
\emph{ICML 2010}.

\bibitem{settles2009}
Settles, B. (2009).
Active learning literature survey.
\emph{CS Technical Report 1648}, UW-Madison.

\bibitem{solomonoff1964}
Solomonoff, R. J. (1964).
A formal theory of inductive inference. Part I.
\emph{Information and Control}, 7(1), 1--22.

\bibitem{borkar2008}
Borkar, V. S. (2008).
\emph{Stochastic Approximation: A Dynamical Systems Viewpoint}. Cambridge University Press.

\bibitem{kushner2003}
Kushner, H. J. \& Yin, G. G. (2003).
\emph{Stochastic Approximation and Recursive Algorithms and Applications} (2nd ed.). Springer.

\bibitem{polyak1992}
Polyak, B. T. \& Juditsky, A. B. (1992).
Acceleration of stochastic approximation by averaging.
\emph{SIAM Journal on Control and Optimization}, 30(4), 838--855.

\bibitem{khalil2002}
Khalil, H. K. (2002).
\emph{Nonlinear Systems} (3rd ed.). Prentice Hall.

\bibitem{hazan2016}
Hazan, E. (2016).
Introduction to online convex optimization.
\emph{Foundations and Trends in Optimization}, 2(3-4), 157--325.

\bibitem{joulani2016}
Joulani, P. et al. (2016).
Online learning under delayed feedback.
\emph{JMLR}, 17(1), 1--54.

\bibitem{kleijn2012}
Kleijn, B. J. K. \& van der Vaart, A. W. (2012).
The Bernstein-von Mises theorem under misspecification.
\emph{Electronic Journal of Statistics}, 6, 354--381.

\bibitem{ghosh2003}
Ghosh, J. K. \& Ramamoorthi, R. V. (2003).
\emph{Bayesian Nonparametrics}. Springer.

\bibitem{robbins1971}
Robbins, H. \& Siegmund, D. (1971).
A convergence theorem for nonnegative almost supermartingales.
\emph{Optimization Methods in Statistics}, 233--257.

\bibitem{bach2011}
Bach, F. \& Moulines, E. (2011).
Non-asymptotic analysis of stochastic approximation algorithms for machine learning.
\emph{NIPS 2011}.

\bibitem{karimi2016}
Karimi, H., Nutini, J. \& Schmidt, M. (2016).
Linear convergence of gradient and proximal-gradient methods under the Polyak-{\L}ojasiewicz condition.
\emph{ECML PKDD 2016}.

\bibitem{benaim1999}
Bena\"im, M. (1999).
Dynamics of stochastic approximation algorithms.
\emph{S\'eminaire de Probabilit\'es}, 1709, 1--68.

\bibitem{pollard1981}
Pollard, D. (1981).
Strong consistency of k-means clustering.
\emph{Annals of Statistics}, 9(1), 135--140.

\bibitem{ghadimi2013}
Ghadimi, S. \& Lan, G. (2013).
Stochastic first- and zeroth-order methods for nonconvex stochastic programming.
\emph{SIAM Journal on Optimization}, 23(4), 2341--2368.

\bibitem{agarwal2012}
Agarwal, A. et al. (2012).
Information-theoretic lower bounds on the oracle complexity of stochastic convex optimization.
\emph{IEEE Transactions on Information Theory}, 58(5), 3235--3249.

\bibitem{bendavid2010}
Ben-David, S. et al. (2010).
A theory of learning from different domains.
\emph{Machine Learning}, 79(1), 151--175.

\bibitem{david1985}
David, P. A. (1985).
Clio and the Economics of QWERTY.
\emph{American Economic Review}, 75(2), 332--337.

\bibitem{arthur1989}
Arthur, W. B. (1989).
Competing technologies, increasing returns, and lock-in by historical events.
\emph{Economic Journal}, 99(394), 116--131.

\bibitem{strogatz2018}
Strogatz, S. H. (2018).
\emph{Nonlinear Dynamics and Chaos} (2nd ed.). CRC Press.

\bibitem{tsybakov2009}
Tsybakov, A. B. (2009).
\emph{Introduction to Nonparametric Estimation}. Springer.

\bibitem{natarajan2013}
Natarajan, N. et al. (2013).
Learning with noisy labels.
\emph{NIPS 2013}.

\end{thebibliography}

\end{document}

\bibitem{he2016}
He, K., Zhang, X., Ren, S., \& Sun, J. (2016).
Deep residual learning for image recognition.
\emph{Proceedings of the IEEE Conference on Computer Vision and Pattern Recognition (CVPR)}, 770--778.
